\documentclass[11pt]{article}
\usepackage[margin=1in]{geometry}
\usepackage{amsmath,amssymb,amsthm}
\usepackage{parskip}
\usepackage{booktabs}

\newtheorem{theorem}{Theorem}
\newtheorem{lemma}[theorem]{Lemma}
\newtheorem{corollary}[theorem]{Corollary}
\newtheorem{proposition}[theorem]{Proposition}
\theoremstyle{definition}
\newtheorem{definition}[theorem]{Definition}
\theoremstyle{remark}
\newtheorem{remark}[theorem]{Remark}

\newcommand{\Z}{\mathbb{Z}}
\newcommand{\orb}[1]{\{#1\}}

\title{Disproof of conjecture \texttt{00000000153}}
\author{earthking11}
\date{2026-09-13}

\begin{document}
\maketitle

\section*{Verdict: FALSE}

We disprove conjecture \texttt{00000000153} as stated. The refutation is
unconditional and elementary: the finite discrete system
$X = \Z/4 = \{0,1,2,3\}$ with $Tx = x+1$ is minimal, yet for $x = 0$ the
prime-index orbit $\{T^{p_n}x : n \ge 1\}$ is $\{1,2,3\}$, whose closure in the
discrete topology is itself and is therefore not all of $X$. In fact the
conclusion fails for \emph{every} $x \in X$.

\section{The conjecture}

Quoted verbatim from \texttt{conjectures/00000000153.md}:

\begin{quote}
\textbf{English.} Conjecture: Minimal transitivity of prime orbits: in any
minimal topological dynamical system $(X, T)$, for every $x \in X$ the orbit
closure $\{T^{p_n}x : n \ge 1\}$ equals $X$.
\end{quote}

\noindent The Chinese statement filed with the conjecture is reproduced
verbatim in \texttt{README.md} (together with this English text). It is not
typeset here, so that \texttt{main.tex} has no CJK font requirement and
compiles with a plain \texttt{tectonic main.tex}. The two formulations are
translations of one another, and both leave the quantifier over $x$ and the
class of systems unrestricted; the refutation therefore applies to either
reading.

The statement leaves the following definitions implicit; we fix the standard
ones, so that the refutation cannot be dismissed as a misreading.

\begin{definition}[Topological dynamical system]
A topological dynamical system is a topological space $X$ together with a
continuous map $T : X \to X$. (In the standard setting of topological dynamics
one usually assumes $X$ compact Hausdorff; the witness below is compact
Hausdorff, so it satisfies that convention.)
\end{definition}

\begin{definition}[Minimal system]
$(X,T)$ is \emph{minimal} if every orbit is dense in $X$; equivalently, the
only closed $T$-invariant subset of $X$ is $X$ itself. In a minimal system the
full orbit $\{T^n x : n \ge 0\}$ is dense for every $x \in X$.
\end{definition}

\begin{definition}[Orbit closure and the prime sequence]
$\overline{A}$ denotes the closure of $A \subseteq X$. The sequence
$(p_n)_{n\ge 1} = (2, 3, 5, 7, 11, \dots)$ is the sequence of primes in
increasing order, so $p_1 = 2$ and $p_n$ is the $n$-th prime. The
\emph{prime-index orbit} of $x$ is the set
\[
P(x) \;=\; \{T^{p_n}x : n \ge 1\} \;=\; \{T^{p}x : p \text{ prime}\},
\]
the two descriptions agreeing because $n \mapsto p_n$ is a bijection onto the
set of primes. The conjecture asserts $\overline{P(x)} = X$ for every $x$.
\end{definition}

\section{The witness: $X = \Z/4$ with the shift}

\begin{definition}[The system]
Let $X = \Z/4 = \{0,1,2,3\}$ with the discrete topology, and let
\[
T : X \to X, \qquad T x = x + 1 \pmod 4 .
\]
$X$ is finite, hence compact, and discrete, hence Hausdorff and metric: it is
a compact metric space. $T$ is continuous (every map on a discrete space is),
so $(X,T)$ is a topological dynamical system in the strictest sense.
\end{definition}

The dynamics is a single cycle:
\[
0 \;\xrightarrow{\;T\;}\; 1 \;\xrightarrow{\;T\;}\; 2
\;\xrightarrow{\;T\;}\; 3 \;\xrightarrow{\;T\;}\; 0 .
\]

\begin{lemma}[$T$ is a $4$-cycle, hence $(X,T)$ is minimal]
\label{lem:minimal}
For every $x \in X$ the full orbit $\{T^n x : n \ge 0\}$ equals $X$.
Consequently $(X,T)$ is minimal.
\end{lemma}

\begin{proof}
Iterating, $T^n x = x + n \pmod 4$ for all $n \ge 0$. For fixed $x$, as $n$
runs over $\{0,1,2,3\}$ the values $x+n \bmod 4$ run over all four residues
$x, x+1, x+2, x+3$, i.e.\ over $X$; and $T^4 x = x$. Hence
$\{T^n x : n \ge 0\} = X$. Every orbit is therefore dense (indeed equal to
$X$), which is minimality.
\end{proof}

\begin{remark}
The full orbit of every point is \emph{exactly} $X$, not merely dense. So the
alternative reading of the conjecture that replaces $P(x)$ by the full orbit
$\{T^{p_n}x\}$ with ``$p_n = n$'' is trivially true here; but that is a
different, degenerate statement (see Section~\ref{sec:caveat}).
\end{remark}

\section{No prime is divisible by $4$}

\begin{lemma}[$4 \nmid p$ for every prime $p$]
\label{lem:four}
If $p$ is prime, then $4 \nmid p$.
\end{lemma}

\begin{proof}
Suppose $4 \mid p$. Then $p = 4k$ for some $k \in \Z$, so $p$ is even, i.e.\
$2 \mid p$. Since $p$ is prime and $2 \mid p$, the divisor $2$ of $p$ must
satisfy $2 = 1$ or $2 = p$; the first is false, so $p = 2$. But then
$4 \mid 2$ would be required, which is false. Hence $4 \nmid p$.
\end{proof}

Equivalently: the only even prime is $2$, and $4 \nmid 2$. Therefore for every
prime $p$,
\[
p \bmod 4 \;\in\; \{1, 2, 3\},
\]
and the value $0$ is never attained. The value set is exactly $\{1,2,3\}$,
witnessed by
\[
2 \bmod 4 = 2, \qquad 3 \bmod 4 = 3, \qquad 5 \bmod 4 = 1 .
\]

\section{The prime-index orbit omits its base point}

\begin{theorem}[Prime orbit of $x$]
\label{thm:prime}
For every $x \in X$,
\[
P(x) \;=\; \{T^p x : p \text{ prime}\} \;=\; X \setminus \{x\}.
\]
In particular $P(0) = \{1,2,3\}$, and $\overline{P(0)} = \{1,2,3\} \neq X$ in
the discrete topology.
\end{theorem}

\begin{proof}
Since $T^n x = x + n \pmod 4$, for a prime $p$ we have
\[
T^p x \;=\; x + p \pmod 4 \;=\; x + (p \bmod 4) \pmod 4 .
\]
By Lemma~\ref{lem:four}, $p \bmod 4 \in \{1,2,3\}$, so
$T^p x \neq x \pmod 4$: the base point $x$ is never in $P(x)$.

Conversely, fix $y \in X$ with $y \neq x$, and let $d = y - x \bmod 4 \in
\{1,2,3\}$. Choose a prime $p \equiv d \pmod 4$:
\[
d = 1 \Rightarrow p = 5, \qquad d = 2 \Rightarrow p = 2, \qquad
d = 3 \Rightarrow p = 3 .
\]
Then $T^p x = x + d = y \pmod 4$, so $y \in P(x)$. Hence
$P(x) = X \setminus \{x\}$.

For $x = 0$ this is $P(0) = \{1,2,3\}$. Because $X$ is discrete, every subset
is closed, so $\overline{P(0)} = P(0) = \{1,2,3\}$, which omits $0$ and is a
proper subset of $X = \{0,1,2,3\}$.
\end{proof}

\begin{corollary}[Refutation]
\label{cor:refute}
$(X,T)$ is a minimal topological dynamical system, but the conjecture's
conclusion
\[
\forall x \in X,\quad \overline{\{T^{p_n}x : n \ge 1\}} = X
\]
is false. Indeed, for every $x \in X$,
\[
\overline{P(x)} \;=\; X \setminus \{x\} \;\subsetneq\; X ,
\]
so the conclusion fails at every single point $x$; a fortiori it fails at the
point $x = 0$ exhibited in the conjecture's own quantifier.
\end{corollary}

\begin{proof}
Minimality is Lemma~\ref{lem:minimal}. The closure computation is
Theorem~\ref{thm:prime} together with discreteness. A proper subset of $X$ is
not equal to $X$, so the asserted equality fails for each $x$.
\end{proof}

\section{Caveat: scope of the refutation}
\label{sec:caveat}

\noindent\textbf{This caveat is part of the result and should be read
together with it.}

\begin{enumerate}
\item \textbf{The witness has isolated points, and is finite.} $X = \Z/4$ is a
finite discrete space, so every point is isolated. Some authors in topological
dynamics study only minimal systems on compact metric spaces \emph{without
isolated points} (or simply \emph{infinite} spaces). $X = \Z/4$ falls outside
that additional, \emph{unstated} hypothesis. If one adds such a hypothesis ---
``$X$ is a compact metric space with no isolated points'', or just ``$X$ is
infinite'' --- then the conjecture as modified is not touched by this
counterexample, and its status remains open here.

\item \textbf{But the file states no such hypothesis.} The statement in
\texttt{conjectures/00000000153.md} says ``in any minimal topological dynamical
system $(X,T)$'', with no cardinality or perfectness restriction, and the
quantification is explicitly ``for every $x \in X$''. Under the literal
statement, $(\Z/4, +1)$ is an admissible minimal system and $x = 0$ is an
admissible point, so the literal statement is refuted. The finite case is not
a technicality outside the intended domain unless the domain is restricted in
writing; it is a minimal topological dynamical system in the standard sense.

\item \textbf{If instead ``orbit closure'' is read as the closure of the full
orbit, the conjecture collapses rather than being rescued.} In a minimal system
the full orbit $\{T^n x : n \ge 0\}$ is dense by definition, so
$\overline{\{T^n x : n \ge 0\}} = X$ holds for every $x$ \emph{by definition}
of minimality. Reading the ambiguous notation $\{T^{p_n}x : n \ge 1\}$ as
the full orbit (i.e.\ $p_n = n$) makes the assertion a tautology for every
minimal system. That is a collapse of the statement to its definition, not a
substantive theorem and not a repair of the prime-index claim: the file's
notation $p_n$ is the standard one for the $n$-th prime, and the conjecture is
titled ``prime orbits''.
\end{enumerate}

In short: the literal conjecture is false, the refutation is elementary and
rigorous, and the only way to avoid it is to add a hypothesis (infiniteness /
no isolated points) that the file does not state. The sharpness note is that
under that added hypothesis the witness is excluded and the conjecture is
\emph{untouched} --- which is precisely why the statement as written is
defective.

\begin{remark}[On the role of $p_1 = 2$]
The prime sequence starts at $p_1 = 2$, and $p = 2$ is included in $P(x)$. This
does not change the conclusion: $2 \bmod 4 = 2 \neq 0$, so even the even prime
contributes a nonzero residue and $P(x)$ still omits $x$.
\end{remark}

\section{Formalisation}

The refutation is formalised in the Lean~4 project under \texttt{lean4/}
(core Lean only --- no Mathlib, no \texttt{sorry}). The system is
$X = \texttt{Fin}\ 4$, the shift is $\texttt{T}\ x = x+1$, the iterate is
\texttt{iter}, and primality is defined directly because core Lean has no
\texttt{Nat.Prime}. The main formal statements are

\begin{quote}\ttfamily
theorem minimal (x y : Fin 4) : FullOrbit x y\\
theorem four\_not\_dvd\_of\_isPrime (hp : IsPrime p) : not (4 | p)\\
theorem primeOrbitAt\_not\_self (x : Fin 4) : not (PrimeOrbitAt x x)\\
theorem primeOrbitClosure\_not\_univ (x : Fin 4) :\\
\hspace*{2em}not (forall y : Fin 4, PrimeOrbitClosure x y)\\
theorem conjecture\_00000000153\_false :\\
\hspace*{2em}IsMinimal and forall x : Fin 4, not (forall y : Fin 4, PrimeOrbitClosure x y)
\end{quote}

\noindent Here \texttt{FullOrbit} is the full-orbit reachability relation (its
totality is minimality), \texttt{PrimeOrbitAt} is the prime-index orbit, and
\texttt{PrimeOrbitClosure} is the prime-index orbit with the closure modelled
by the set itself, which is correct because $X$ is finite and discrete. The
project also has \texttt{conjecture\_00000000153\_false\_at\_zero}, the same
statement specialised to $x = 0$. The axiom audit
(\texttt{lake env lean Check.lean}) reports only \texttt{propext} and
\texttt{Quot.sound} for every theorem, and no \texttt{sorryAx}. See
\texttt{lean4/README.md} for the full statement table and scope note.

\section{Reproducibility}

\begin{quote}\ttfamily
\$ python3 reproduce.py\\
... (enumerates all primes p < 1e6; checks the residue set,\\
... the 4-cycle, minimality, and each prime orbit)\\
PASS: all checks verified; conjecture 00000000153 is FALSE.
\end{quote}

\texttt{reproduce.py} uses only the Python~3 standard library. It sieves all
$78498$ primes $p < 10^6$, verifies that no such prime is divisible by $4$ and
that the residue set $\{p \bmod 4\}$ is exactly $\{1,2,3\}$; verifies that $T$
is a $4$-cycle and that the full orbit of every $x$ is $X$ (minimality); and
verifies that for each $x \in \{0,1,2,3\}$ the prime-index orbit is exactly
$X \setminus \{x\}$, so its closure is a proper subset of $X$. It prints
\texttt{PASS} and exits $0$ exactly when every check holds, and exits non-zero
otherwise.

The Lean project is built and audited by

\begin{quote}\ttfamily
cd lean4\\
lake build\\
lake env lean Check.lean
\end{quote}

The document is compiled by

\begin{quote}\ttfamily
tectonic main.tex\\
(equivalently: tectonic --outdir build main.tex)
\end{quote}

\noindent producing \texttt{build/main.pdf} (with \texttt{tectonic main.tex}
the file is written to \texttt{main.pdf} in the submission root and then placed
at \texttt{build/main.pdf}, the same location as the template submission).

\section{Numerical vantage point}

\begin{center}
\begin{tabular}{r l l l c}
\toprule
$x$ & full orbit $\{T^n x\}$ & prime orbit $P(x)$ & omitted & $\overline{P(x)} = X$? \\
\midrule
$0$ & $\{0,1,2,3\}$ & $\{1,2,3\}$     & $0$ & no \\
$1$ & $\{0,1,2,3\}$ & $\{0,2,3\}$     & $1$ & no \\
$2$ & $\{0,1,2,3\}$ & $\{0,1,3\}$     & $2$ & no \\
$3$ & $\{0,1,2,3\}$ & $\{0,1,2\}$     & $3$ & no \\
\bottomrule
\end{tabular}
\end{center}

Every full orbit is $X$ (the system is minimal), yet every prime-index orbit
omits exactly one point (its base point), so no prime-index orbit closure equals
$X$. The conjecture fails at all four points of the minimal system.

\section{Status against the submission rules}

\begin{itemize}
\item \textbf{LaTeX source} --- present (\texttt{main.tex}), a standalone
\texttt{article} using only \texttt{geometry}, \texttt{amsmath},
\texttt{amssymb}, \texttt{amsthm}, \texttt{parskip}, and \texttt{booktabs}; it
compiles with \texttt{tectonic main.tex}.
\item \textbf{PDF document} --- present at \texttt{build/main.pdf}, produced by
\texttt{tectonic main.tex}.
\item \textbf{Lean 4 project} --- present under \texttt{lean4/}, with
\texttt{lean-toolchain} (pinned to \texttt{leanprover/lean4:v4.33.1}),
\texttt{lakefile.toml} (library \texttt{Main}, project \texttt{tlmc153}),
\texttt{Main.lean} (core Lean only, no \texttt{sorry}), \texttt{Check.lean}
(\texttt{\#print axioms} audit), and \texttt{README.md}. The audit reports only
\texttt{[propext, Quot.sound]} and no \texttt{sorryAx}.
\item \textbf{Reproduction} --- \texttt{python3 reproduce.py} passes with exit
code $0$ and prints \texttt{PASS}.
\end{itemize}

\end{document}
